\documentclass[uplatex, dvipdfmx, a4paper, 11pt]{jsarticle}

\usepackage[top=25truemm,bottom=25truemm,left=25truemm,right=25truemm]{geometry}
\usepackage{amsmath, amssymb, mathtools, bm}
\usepackage[dvipdfmx]{graphicx}
\usepackage{float}
\usepackage{booktabs}
\usepackage{caption}
\usepackage[dvipdfmx, bookmarks=true, setpagesize=false, colorlinks=true, linkcolor=black, urlcolor=blue, citecolor=black]{hyperref}
\usepackage{pxjahyper}
\usepackage{listings} 
\usepackage{xcolor}
\usepackage{algorithm}
\usepackage{algorithmicx}
\usepackage{algpseudocode}
\usepackage{tikz}

\definecolor{codegreen}{rgb}{0,0.6,0}
\definecolor{codegray}{rgb}{0.5,0.5,0.5}
\definecolor{codepurple}{rgb}{0.58,0,0.82}
\definecolor{backcolour}{rgb}{0.95,0.95,0.92}
\lstdefinestyle{mystyle}{
    backgroundcolor=\color{backcolour},   
    commentstyle=\color{codegreen},
    keywordstyle=\color{magenta},
    numberstyle=\tiny\color{codegray},
    stringstyle=\color{codepurple},
    basicstyle=\ttfamily\footnotesize,
    breakatwhitespace=false,         
    breaklines=true,                 
    captionpos=b,                    
    keepspaces=true,                 
    numbers=left,                    
    numbersep=5pt,                  
    showspaces=false,                
    showstringspaces=false,
    showtabs=false,                  
    tabsize=2,
    frame=single
}
\lstset{style=mystyle}

\usetikzlibrary{positioning, arrows.meta, calc}

\title{環境情報学 2026年度 第1回授業後課題\\ \large ー葉っぱを探そうの会ー}
\author{
    慶應義塾大学 総合政策学部 / 環境情報学部 \\[3mm]
    \begin{tabular}{ll}
        学籍番号: 72603857 & 氏名: 鈴木 真理 \\
        学籍番号: 72646131 & 氏名: 長谷部 舞衣 \\
        学籍番号: 72600825 & 氏名: 上木 久里 \\
    \end{tabular}
}
\date{\today}

\begin{document}

\maketitle

\begin{abstract}
    与えられた葉っぱをChatGPTの画像認識，Geminiの画像認識を用いて「\textbf{スギ (\textit{Cryptomeria japonica})}」であると特定することができた．
\end{abstract}

\section{葉っぱの特徴}


\section{葉っぱの特定方法}
開始当初は，葉っぱの特徴をもとに自然言語で検索する方法を試みた．しかし，葉っぱの特徴を正確に表現することが難しく，検索結果も不十分であった．そこで，画像認識技術を活用することにした．ChatGPTの画像認識機能を使用して，葉っぱの画像をアップロードし，特徴を分析したところ，「スギ (\textit{Cryptomeria japonica})」であると特定された．さらに，Geminiの画像認識機能を使用して同様の分析を行ったところ，同じく「スギ (\textit{Cryptomeria japonica})」であると特定された．最後に福沢育林の会が公開する『慶應義塾の環境・安全・健康への取り組み2008紹介』という資料を参照し，スギが慶應義塾大学のキャンパス内に植えられていることを確認した．これらの結果から，与えられた葉っぱは「\textbf{スギ (\textit{Cryptomeria japonica})}」であると結論づけることができた．

\section{植生場所の発見手順}
葉っぱが「\textbf{スギ (\textit{Cryptomeria japonica})}」であることが特定できたため，植生場所を特定するために以下の手順を踏んだ．まず，スギがどのような環境で育つかを調査した．スギは日本の代表的な針葉樹であり，温暖湿潤な気候を好むことがわかった．次に，慶應義塾大学のキャンパス内でスギが植えられている場所を特定するために，キャンパスマップや植生図を参照した．さらに，キャンパス内を実際に歩いて確認し、スギが植えられている場所を特定した．最終的に，スギが植えられている場所はキャンパスの北側の体育館裏であることがわかった．

\section{植生場所と考察}


\section{得た学び}
この課題を通して，物事を正確に表現/伝達することの難しさと重要性を学んだ．

\end{document}